% Appendix Table A.6: whole-task versus subtask-based scores. Generated by code/tables/tabB-whole-task.py
% from the whole-task comparison files (data/scoring/whole_task_*.csv). Do not edit by hand.
\begin{table}[htbp]
\centering
\begin{threeparttable}
\caption{Scoring the task as a whole versus scoring its subtasks}
\label{tab:wholetask}
\footnotesize\setlength{\tabcolsep}{5pt}
\begin{spacing}{1}
\begin{tabular}{@{}lrrrr@{}}
\toprule
 & Claude Haiku 4.5 & GPT-5.6-luna & Gemini 3.1 Flash-Lite & Three-model average \\
\midrule
\multicolumn{5}{@{}l}{\emph{Mean exposure score at stage 4}} \\
Subtask-based & 47.1 & 40.9 & 42.7 & -- \\
Whole task & 32.6 & 38.9 & 33.7 & -- \\
\addlinespace
\multicolumn{5}{@{}l}{\emph{Rank correlation between the two methods, same model}} \\
Stage-4 level & 0.84 & 0.86 & 0.87 & -- \\
Change, stage 1 to 4 & 0.77 & 0.78 & 0.79 & -- \\
\addlinespace
\multicolumn{5}{@{}l}{\emph{Distinct values of the time-saved rating at stage 4}} \\
Subtask-based & 726 & 863 & 682 & -- \\
Whole task & 44 & 100 & 26 & -- \\
\addlinespace
\multicolumn{5}{@{}l}{\emph{Rank correlation of the stage-1 to stage-4 change with the other two models, average}} \\
Subtask-based & 0.76 & 0.78 & 0.80 & -- \\
Whole task & 0.84 & 0.84 & 0.86 & -- \\
\addlinespace
\multicolumn{5}{@{}l}{\emph{Correlation with adoption across 141 occupations, change from stage 1 to 3}} \\
Subtask-based & 0.74 & 0.72 & 0.72 & 0.74 \\
Whole task & 0.69 & 0.75 & 0.64 & 0.70 \\
\addlinespace
\multicolumn{5}{@{}l}{\emph{Correlation with adoption across 141 occupations, change from stage 1 to 4}} \\
Subtask-based & 0.67 & 0.75 & 0.71 & 0.71 \\
Whole task & 0.66 & 0.76 & 0.65 & 0.70 \\
\addlinespace
\multicolumn{5}{@{}l}{\emph{Test-retest rank correlation of the time-saved rating at stage 4, 1,000-task pilot}} \\
Subtask-based & -- & -- & 0.97 & -- \\
Whole task & -- & -- & 0.99 & -- \\
\bottomrule
\end{tabular}
\end{spacing}
\begin{tablenotes}\notesize
\item Notes: Each model scored all 19,281 O*NET tasks twice: once subtask by subtask with the fixed decomposition and the three prompts of Section 2, and once as a whole, with the task statement alone and the same three questions asked as percentages of the task, using the same stage definitions, typology, and monotonicity rules. Whole-task scores were obtained through OpenRouter on October 8, 2026, at temperature zero. The composite is the same average of the three ratings in both cases. Correlations are across tasks. The adoption rows aggregate each score to census occupations as in Section 2 and correlate the change with the share of workers using generative AI at work in the Real-Time Population Survey, weighted by respondents, for the 141 occupations with at least 20 respondents. The test-retest rows repeat the 1,000-task pilot sample a second time with the same settings for one model.
\end{tablenotes}
\end{threeparttable}
\end{table}
